% A fully connected GAN is shown as one teaching implementation, not a restriction.
\begin{tikzpicture}[x=1mm,y=1mm, font=\fontsize{8}{9.3}\selectfont,
  line width=.65pt, draw=FigureStroke, text=FigureInk,
  arr/.style={-{Stealth[length=1.65mm,width=1.05mm]}},
  layer/.style={rounded corners=1.8pt,draw=FigureBlue,fill=FigureBlueFill,
    minimum width=18mm,minimum height=14mm,align=center,inner sep=1.1mm},
  dlayer/.style={layer,draw=FigureGreen,fill=FigureGreenFill,minimum width=22mm,minimum height=12mm}]
  \path[use as bounding box] (0,0) rectangle (150,90);
  \draw[FigureBlue,dashed,rounded corners=1.8pt] (2,42) rectangle (104,85);
  \node[anchor=west,text=FigureBlue] at (3,87) {生成路径：采样时保留};
  \node[text=FigureBlue] at (49,79) {生成器 $G_{\symbf{\theta}}$};
  \node[layer,minimum width=12mm,minimum height=14mm] (z) at (10,61) {$\symbf z$};
  \node[layer] (g1) at (30,61) {全连接\\激活};
  \node[layer] (g2) at (53,61) {全连接\\激活};
  \node[layer,minimum width=17mm] (g3) at (75,61) {观测层\\重排};
  \draw[arr] (z.east)--(g1.west);
  \draw[arr] (g1.east)--(g2.west);
  \draw[arr] (g2.east)--(g3.west);
  \draw[arr] (g3.east)--(90,61);
  \node[anchor=north] at (10,51) {随机输入};
  \node[anchor=north] at (96,51) {生成样本 $\tilde{\symbf x}$};
  \draw[FigureBlue,fill=FigureBlueFill] (90,55) rectangle (102,67);
  \foreach \i in {1,2,3} {
    \draw[FigureBlue!55,line width=.3pt] (90+3*\i,55)--(90+3*\i,67);
    \draw[FigureBlue!55,line width=.3pt] (90,55+3*\i)--(102,55+3*\i);
  }
  \fill[FigureBlue] (93,58) rectangle (96,64);
  \fill[FigureBlue] (96,61) rectangle (99,64);
  \draw[FigureGreen,dashed,rounded corners=1.8pt] (115,20) rectangle (148,85);
  \node[text=FigureGreen] at (131.5,87) {共享判别器 $D_{\symbf\phi}$};
  \node[dlayer] (d1) at (131.5,67) {隐藏层\\激活};
  \node[dlayer] (d2) at (131.5,49) {线性层\\单一得分};
  \node[dlayer] (d3) at (131.5,31) {sigmoid};
  \draw[arr] (102,61)--(110,61)|-(d1.west);
  \draw[arr] (d1.south)--(d2.north);
  \draw[arr] (d2.south)--(d3.north);
  \draw[arr] (d3.south)--(131.5,13);
  \node[anchor=north] at (131.5,11) {$D_{\symbf\phi}(\symbf x)\in(0,1)$};
  \draw[FigureGreen,fill=FigureGreenFill] (90,21) rectangle (102,33);
  \foreach \i in {1,2,3} {
    \draw[FigureGreen!55,line width=.3pt] (90+3*\i,21)--(90+3*\i,33);
    \draw[FigureGreen!55,line width=.3pt] (90,21+3*\i)--(102,21+3*\i);
  }
  \fill[FigureGreen] (93,24) rectangle (96,30);
  \fill[FigureGreen] (96,27) rectangle (99,30);
  \draw[arr] (102,27)--(112,27)--(112,64)--(120.5,64);
  \node[anchor=north] at (96,18) {真实样本 $\symbf x$};
  \node[anchor=west] at (3,29) {来源监督：真实为 $1$，生成为 $0$};
  \node[anchor=west,font=\fontsize{7.5}{9}\selectfont\figurenote] at (3,22) {训练 $G$ 时，经过 $D$ 回传输入梯度};
  \node[anchor=west,font=\fontsize{7.5}{9}\selectfont\figurenote] at (3,15) {此时固定 $\symbf\phi$，只更新 $\symbf\theta$};
\end{tikzpicture}
